% =====================================================================
%  TA-IND-01 - Ingenieria de Requerimientos (ISR-401)
%  Facultad de Ciencias de la Computacion - UTEQ
%  Estudiante: Roselyn Andreina Sanchez Centeno
%  Sistema analizado: aplicacion movil y web Shein
%
%  ARCHIVO PRINCIPAL. Compilar con pdfLaTeX + BibTeX:
%      pdflatex main.tex
%      bibtex main
%      pdflatex main.tex
%      pdflatex main.tex
%  Alternativa: tectonic main.tex
%
%  ANTES DE ENTREGAR: reemplazar cada marca roja [FALTA: ...] por el
%  texto definitivo y comprobar que no queda ninguna en main.pdf.
% =====================================================================
\documentclass[12pt,a4paper]{article}

\usepackage[T1]{fontenc}
\usepackage[utf8]{inputenc}
\usepackage{lmodern}
\usepackage[spanish,es-noshorthands]{babel}
\usepackage[margin=2.5cm]{geometry}
\usepackage{microtype}
\usepackage{graphicx}
\usepackage{booktabs}
\usepackage{longtable}
\usepackage{array}
\usepackage{tabularx}
\usepackage{pdflscape}
\usepackage{enumitem}
\usepackage{xcolor}
\usepackage{caption}
\usepackage{cite}
\usepackage[hidelinks]{hyperref}

\graphicspath{{figuras/}}

% ---------------------------------------------------------------------
% Ayudas de trabajo (no forman parte del contenido entregable)
% ---------------------------------------------------------------------
\newcommand{\falta}[1]{\textcolor{red}{\textbf{[FALTA: #1]}}}

\newcommand{\captura}[3]{%
  \begin{figure}[htbp]
    \centering
    \IfFileExists{figuras/#1}{%
      \includegraphics[width=0.58\linewidth]{#1}%
    }{%
      \fbox{\parbox[c][5.0cm][c]{0.58\linewidth}{%
        \centering \falta{captura \texttt{#1}}}}%
    }%
    \caption{#2}
    \label{#3}
  \end{figure}%
}

\captionsetup{font=small,labelfont=bf}
\setlength{\parindent}{0pt}
\setlength{\parskip}{0.45em}
\renewcommand{\arraystretch}{1.2}

% ---------------------------------------------------------------------
\begin{document}

% ============================ PORTADA ================================
\begin{titlepage}
  \centering
  \vspace*{1.4cm}
  {\Large\bfseries Requisitos de un sistema conocido a partir de\\[4pt]
  ISO/IEC/IEEE 29148:2018\par}
  \vspace{1.0cm}
  {\large TA-IND-01 --- Trabajo autonomo individual\par}
  \vspace{0.4cm}
  {\large Ingenieria de Requerimientos (ISR-401) --- Cuarto nivel --- Paralelo A\par}
  \vspace{0.4cm}
  {\large Facultad de Ciencias de la Computacion\par}
  {\large Universidad Tecnica Estatal de Quevedo\par}
  \vspace{1.4cm}
  {\large\bfseries Roselyn Andreina Sanchez Centeno\par}
  \vspace{0.3cm}
  {\large Sistema analizado: Shein (aplicacion movil y sitio web)\par}
  \vspace{1.4cm}
  {\large Fecha de cierre: sabado 7 de noviembre de 2026, 23:59\par}
  \vfill
\end{titlepage}

% =====================================================================
\section{Descripcion del sistema}
\label{sec:sistema}

Shein es una plataforma para realizar compras en linea que vende, en su
mayoria, ropa. Esta dividida por categorias como el genero, las tallas,
los colores y los precios. La plataforma permite filtrar por categoria,
realizar pedidos, pagar mediante tarjeta de debito o credito y dar
seguimiento al pedido. La usan todo tipo de personas, desde el celular
o desde la pagina web.

Las Figuras~\ref{fig:cap1} a~\ref{fig:cap4} son capturas propias de la
aplicacion y respaldan los requisitos de la Seccion~\ref{sec:requisitos}.

\captura{cap01_pantalla_inicio.jpeg}{Pantalla de inicio de la aplicacion
Shein.}{fig:cap1}

\captura{cap02_busqueda_catalogo.jpeg}{Busqueda y filtros del catalogo de
productos.}{fig:cap2}

\captura{cap03_carrito_compra.jpeg}{Carrito de compra con los articulos
seleccionados.}{fig:cap3}

\captura{cap04_pantalla_pago.jpeg}{Pantalla de pago y seleccion del metodo
de envio.}{fig:cap4}

% =====================================================================
\section{Resumen analitico de ISO/IEC/IEEE 29148:2018}
\label{sec:resumen}

Este resumen se basa en las clausulas 1, 3 y 5.2 de
ISO/IEC/IEEE 29148:2018~\cite{iso29148}, complementadas con el modelo de
calidad de producto de ISO/IEC 25010:2023~\cite{iso25010} y con los
fundamentos de ingenieria de requisitos de SWEBOK v4.0~\cite{swebok4}.

\subsection{Que es un requisito bien formado (clausula 5.2.4)}
\falta{explicacion de la construccion de un requisito bien formado segun
la clausula 5.2.4, con cita a~\cite{iso29148}.}

\subsection{Las nueve caracteristicas de un requisito individual (clausula 5.2.5)}
\falta{explicacion de las nueve caracteristicas: necesaria, libre de
implementacion, no ambigua, completa, singular, factible, verificable,
correcta y conforme. Citar la clausula 5.2.5.}

\subsection{Caracteristicas de un conjunto de requisitos (clausula 5.2.6)}
\falta{explicacion de las caracteristicas del conjunto y de por que un
conjunto puede fallar aunque cada requisito sea correcto. Citar la
clausula 5.2.6.}

\subsection{Criterios de lenguaje y terminos prohibidos (clausula 5.2.7)}
\falta{explicacion de los criterios de lenguaje, con ejemplos de terminos
ambiguos que la norma prohibe y su equivalente verificable. Citar la
clausula 5.2.7.}

\subsection{Requisito y atributo (clausula 5.2.8)}
\falta{distincion entre el requisito y sus atributos, y utilidad de los
atributos para la trazabilidad. Citar la clausula 5.2.8.}

% =====================================================================
\section{Glosario}
\label{sec:glosario}

\begin{longtable}{>{\raggedright\arraybackslash}p{3.2cm}
                  >{\raggedright\arraybackslash}p{6.0cm}
                  >{\raggedright\arraybackslash}p{5.0cm}}
\caption{Glosario de diez terminos de la clausula 3 de ISO/IEC/IEEE
29148:2018, con ejemplo en Shein.}
\label{tab:glosario}\\
\toprule
\textbf{Termino} & \textbf{Definicion parafraseada} & \textbf{Ejemplo en Shein} \\
\midrule
\endfirsthead
\toprule
\textbf{Termino} & \textbf{Definicion parafraseada} & \textbf{Ejemplo en Shein} \\
\midrule
\endhead
\bottomrule
\endfoot
\falta{termino 1} & \falta{definicion} & \falta{ejemplo} \\
\falta{termino 2} & \falta{definicion} & \falta{ejemplo} \\
\falta{termino 3} & \falta{definicion} & \falta{ejemplo} \\
\falta{termino 4} & \falta{definicion} & \falta{ejemplo} \\
\falta{termino 5} & \falta{definicion} & \falta{ejemplo} \\
\falta{termino 6} & \falta{definicion} & \falta{ejemplo} \\
\falta{termino 7} & \falta{definicion} & \falta{ejemplo} \\
\falta{termino 8} & \falta{definicion} & \falta{ejemplo} \\
\falta{termino 9} & \falta{definicion} & \falta{ejemplo} \\
\falta{termino 10} & \falta{definicion} & \falta{ejemplo} \\
\end{longtable}

% =====================================================================
\section{Requisitos del sistema}
\label{sec:requisitos}

\begin{longtable}{>{\raggedright\arraybackslash}p{1.6cm}
                  >{\raggedright\arraybackslash}p{12.4cm}}
\caption{Enunciado de los quince requisitos del sistema.}
\label{tab:requisitos}\\
\toprule
\textbf{ID} & \textbf{Enunciado del requisito} \\
\midrule
\endfirsthead
\toprule
\textbf{ID} & \textbf{Enunciado del requisito} \\
\midrule
\endhead
\bottomrule
\endfoot
RF-01 & Cuando el usuario ingrese a Shein, el sistema mostrara el menu
principal de navegacion. \\
RF-02 & Cuando el usuario escriba el nombre o la descripcion de un
producto en la barra de busqueda y pulse buscar, el sistema mostrara los
productos relacionados con su busqueda. \\
RF-03 & Cuando el usuario acceda al menu de navegacion, el sistema
mostrara las categorias principales: ropa de mujer, trajes de bano, talla
grande y ninos. \\
RF-04 & Cuando el usuario seleccione una categoria, el sistema mostrara
los articulos disponibles con sus imagenes, nombres y precios. \\
RF-05 & Cuando el usuario pulse el boton para agregar un articulo al
carrito, el sistema incorporara el producto seleccionado al carrito. \\
RF-06 & Cuando el usuario este en el carrito, el sistema permitira
cambiar la cantidad de cada articulo mediante el selector de unidades. \\
RIE-01 & Cuando el usuario acceda a la seccion de metodos de pago, el
sistema mostrara Visa, Mastercard, Maestro, American Express, Diners
Club, Discover, Apple Pay y UnionPay. \\
RR-01 & El sistema limitara el destino de envio del pedido a Ecuador. \\
RR-02 & El sistema recopilara unicamente los datos personales necesarios
para cumplir con la finalidad del tratamiento, de acuerdo con el
principio de pertinencia y minimizacion establecido en el articulo 10,
literal e, de la LOPDP~\cite{lopdp2021}. \\
RP-01 & Antes de publicar cada version del sistema, el responsable del
proyecto registrara su aprobacion mediante un acta de liberacion. \\
RNF-01 & Cuando el usuario realice una busqueda, el sistema mostrara los
resultados en un tiempo de respuesta cuyo percentil 95 sea de 3 segundos
o menos, bajo una carga de 50 usuarios concurrentes. \\
RNF-02 & Cuando un usuario nuevo utilice el sistema por primera vez,
podra agregar un producto al carrito sin recibir ayuda externa. \\
RNF-03 & Durante el proceso de pago, el sistema protegtera los datos
personales y financieros de los usuarios frente a accesos no
autorizados. \\
RNF-04 & El sistema adaptara la visualizacion de su interfaz a
diferentes tamanos de pantalla, incluyendo celulares, tabletas y
computadoras. \\
RNF-05 & El sistema estara disponible para los usuarios las 24 horas del
dia, durante todos los dias del mes. \\
\end{longtable}

% =====================================================================
\begin{landscape}
\section{Tabla de clasificacion}
\label{sec:clasificacion}

\begin{longtable}{>{\raggedright\arraybackslash}p{1.3cm}
                  >{\raggedright\arraybackslash}p{6.6cm}
                  >{\raggedright\arraybackslash}p{2.4cm}
                  >{\raggedright\arraybackslash}p{4.2cm}
                  >{\raggedright\arraybackslash}p{3.8cm}
                  >{\raggedright\arraybackslash}p{4.2cm}}
\caption{Clasificacion de los requisitos por tipo, caracteristica
ISO/IEC 25010:2023, metrica y criterio de aceptacion.}
\label{tab:clasificacion}\\
\toprule
\textbf{ID} & \textbf{Enunciado (resumido)} & \textbf{Tipo} &
\textbf{ISO/IEC 25010} & \textbf{Metrica} & \textbf{Criterio de aceptacion} \\
\midrule
\endfirsthead
\toprule
\textbf{ID} & \textbf{Enunciado (resumido)} & \textbf{Tipo} &
\textbf{ISO/IEC 25010} & \textbf{Metrica} & \textbf{Criterio de aceptacion} \\
\midrule
\endhead
\bottomrule
\endfoot
RF-01 & Muestra el menu principal de navegacion al ingresar. & Funcional & --- & Cantidad de menus principales visibles. & Se muestra el menu principal al ingresar. \\
RF-02 & Devuelve productos relacionados con el termino buscado. & Funcional & --- & Cantidad de resultados relacionados con la busqueda. & Se muestran productos relacionados con el termino ingresado. \\
RF-03 & Muestra las cuatro categorias principales en el menu. & Funcional & --- & Cantidad de categorias principales visibles. & Se muestran las 4 categorias indicadas. \\
RF-04 & Muestra articulos con imagen, nombre y precio al elegir categoria. & Funcional & --- & Cantidad de articulos con imagen, nombre y precio. & El 100\,\% de los articulos mostrados contiene los tres elementos. \\
RF-05 & Incorpora al carrito el articulo seleccionado. & Funcional & --- & Cantidad de productos agregados correctamente. & El producto aparece en el carrito despues de agregarlo. \\
RF-06 & Permite cambiar la cantidad de cada articulo en el carrito. & Funcional & --- & Cantidad de unidades modificadas correctamente. & La cantidad se actualiza correctamente en el carrito. \\
RIE-01 & Muestra los ocho metodos de pago disponibles. & Interfaz externa & --- & Cantidad de metodos de pago visibles. & Se muestran las 8 opciones de pago identificadas. \\
RR-01 & Limita el destino de envio del pedido a Ecuador. & Restriccion & --- & Cantidad de destinos de envio disponibles. & Ecuador aparece como destino y no se permite otro. \\
RR-02 & Recopila solo los datos personales necesarios (LOPDP). & Restriccion & --- & Cantidad de campos sin finalidad justificada. & El 100\,\% de los campos tiene una finalidad identificada. \\
RP-01 & Exige un acta de liberacion antes de publicar cada version. & Proceso & --- & Cantidad de actas aprobadas por version. & Cada version publicada cuenta con un acta aprobada. \\
RNF-01 & Responde las busquedas en 3 s o menos (p95, 50 usuarios). & Calidad & Eficiencia de desempeno: comportamiento temporal & Tiempo de respuesta en segundos (p95). & En 100 busquedas con 50 usuarios concurrentes, el p95 es de 3 s o menos. \\
RNF-02 & Un usuario nuevo agrega un producto sin ayuda externa. & Calidad & Capacidad de interaccion: capacidad de aprendizaje & Porcentaje de usuarios que completan la tarea en 5 min. & Al menos 9 de 10 usuarios nuevos lo logran en 5 min o menos. \\
RNF-03 & Protege los datos de pago frente a accesos no autorizados. & Calidad & Seguridad: confidencialidad & Numero de accesos no autorizados exitosos. & Cero accesos no autorizados en 20 intentos de prueba. \\
RNF-04 & Adapta la interfaz a celular, tableta y computadora. & Calidad & Flexibilidad: adaptabilidad & Numero de pruebas de visualizacion satisfactorias. & Las 30 pruebas en 375, 768 y 1366 px no presentan desbordamientos. \\
RNF-05 & Permanece disponible 24 horas al dia, todos los dias. & Calidad & Fiabilidad: disponibilidad & Porcentaje de disponibilidad mensual. & Disponibilidad igual o superior al 99\,\%, medida cada 5 min. \\
\end{longtable}
\end{landscape}

% =====================================================================
\section{Autoevaluacion}
\label{sec:autoevaluacion}

\setlength{\tabcolsep}{4pt}
\begin{longtable}{>{\raggedright\arraybackslash}p{1.2cm}
                  >{\raggedright\arraybackslash}p{4.0cm}
                  >{\raggedright\arraybackslash}p{4.6cm}
                  >{\raggedright\arraybackslash}p{4.6cm}}
\caption{Requisitos que fallaron la autoevaluacion: version anterior y
version corregida.}
\label{tab:autoevaluacion}\\
\toprule
\textbf{ID} & \textbf{Caracteristica incumplida} &
\textbf{Version anterior} & \textbf{Version corregida} \\
\midrule
\endfirsthead
\toprule
\textbf{ID} & \textbf{Caracteristica incumplida} &
\textbf{Version anterior} & \textbf{Version corregida} \\
\midrule
\endhead
\bottomrule
\endfoot
\falta{ID} & \falta{5.2.5 / 5.2.7} & \falta{texto} & \falta{texto} \\
\falta{ID} & \falta{5.2.5 / 5.2.7} & \falta{texto} & \falta{texto} \\
\falta{ID} & \falta{5.2.5 / 5.2.7} & \falta{texto} & \falta{texto} \\
\end{longtable}

% =====================================================================
\section{Reflexion}
\label{sec:reflexion}

\falta{reflexion de media pagina sobre el tipo de requisito que costo mas
formular y su causa, con un ejemplo concreto de este trabajo.}

% =====================================================================
\section{Declaracion de uso de inteligencia artificial}
\label{sec:ia}

Se utilizo un asistente de inteligencia artificial como apoyo en la
infraestructura LaTeX del informe: estructura del documento, plantilla de
tablas, archivo de referencias y resolucion de errores de compilacion.
La seleccion del sistema, la descripcion, el analisis de la norma, el
glosario, los requisitos, la autoevaluacion y la reflexion son de autoria
de la estudiante y fueron verificados por ella antes de la entrega.
\falta{completar: nombre y version exacta de la herramienta usada, y como
se verifico lo obtenido.}

% =====================================================================
\bibliographystyle{IEEEtran}
\bibliography{referencias}

\end{document}
